% Process-model stage runtime — \input{} this file.
\begin{table}[H]
\centering

\caption{Computational cost of the process-model stages.}
\label{tab:timing_process_stages}

\vspace{-0.5em}
\footnotesize
\begin{tabular}{p{5.2cm}|r|l|r|r}
\toprule
Stage & Cost & Unit & N & Total (s) \\
\midrule
Process discovery -- Alpha miner & 1.68 & s / process & 6 & 10.1 \\
Process discovery -- Heuristic miner & 5.05 & s / process & 6 & 30.3 \\
Process discovery -- Inductive miner & 5.07 & s / process & 6 & 30.4 \\
Case-duration model (budgeting) & 0.51 & s / process & 6 & 3.0 \\
Case generation -- sampled durations & 1.13 & ms / case & 5,544 & 6.2 \\
Case generation -- ML global durations & 5.37 & ms / case & 5,544 & 29.8 \\
Case generation -- ML per-activity durations & 2.30 & ms / case & 5,544 & 12.8 \\
Process-model scoring & 2.38 & ms / case & 16,632 & 39.5 \\
Complete profile + scoring -- Sensor median & 5.77 & ms / case & 16,632 & 95.9 \\
Complete profile + scoring -- ML Step DTW (proposed) & 436.46 & ms / case & 16,632 & 7,259.2 \\
\bottomrule
\end{tabular}

\vspace{0.5em}

\parbox{\linewidth}{%
\footnotesize
Runtime of the process-model stages, same run. Fitting stages are per process, generation stages per generated case (test split, 6 processes). Case generation is the simulator alone; complete profile adds the curve models over a whole case and the comparison against its real counterpart. Cost is single-process wall clock, a ratio rather than an absolute deployment figure; Total is what this run spent on the stage over its N units, so it scales with the experiment grid (every net crossed with every process).
}

\end{table}
